\section{مقارنة بين محركات السيرفو للتيار المستمر والمتناوب}
\label{sec:ac_dc_comparison}

تاريخياً، كانت محركات التيار المستمر (\textenglish{DC Servo Motors}) هي الخيار السائد في أنظمة التحكم نظراً لسهولة نمذجتها الرياضية وبساطة دوائر التحكم في سرعتها وعزمها. ولكن مع التطور الهائل في إلكترونيات القدرة (\textenglish{Power Electronics})، وتقنيات المعالجات الدقيقة، أصبحت محركات التيار المتناوب (\textenglish{AC Servo Motors}) هي المعيار الذهبي المعتمد في الأتمتة الصناعية الحديثة.

\subsection{خصائص سيرفو التيار المستمر (\textenglish{DC Servo})}
تُستخدم هذه المحركات غالباً في التطبيقات ذات الاستطاعات المنخفضة إلى المتوسطة، أو في الأنظمة التي تعتمد على البطاريات (مثل الروبوتات المتنقلة). وتعتمد إما على المحركات ذات الفرش (\textenglish{Brushed DC}) أو بدون فرش (\textenglish{BLDC}).
\begin{itemize}
    \item \textbf{المزايا:} سهولة التحكم والمناورة، وتكلفة ابتدائية أقل لمشغلات القدرة (\textenglish{Drives}).
    \item \textbf{القيود:} في الأنواع ذات الفرش الكربونية، يتطلب الأمر صيانة دورية بسبب التآكل الميكانيكي للفرش. كما أن كفاءة تبديد الحرارة تكون أقل نظراً لتولد الحرارة في العضو الدوار الداخلي.
\end{itemize}

\subsection{خصائص سيرفو التيار المتناوب (\textenglish{AC Servo})}
تعتمد بشكل أساسي على المحركات التزامنية ذات المغناطيس الدائم (\textenglish{PMSM}). وهي الخيار الأول والتطبيق الفعلي للروبوتات الصناعية وآلات القطع المبرمجة (\textenglish{CNC}).
\begin{itemize}
    \item \textbf{المزايا:} غياب التآكل الميكانيكي للفرش (انعدام الحاجة للصيانة الدورية)، وكفاءة تبريد عالية جداً؛ حيث تقع الملفات النحاسية على العضو الثابت الخارجي، مما يسهل تشتيت الحرارة إلى الهيكل المعدني للآلة. بالإضافة إلى توفير عزم مستقر حتى عند السرعات العالية جداً.
    \item \textbf{القيود:} تتطلب خوارزميات تحكم رياضية معقدة جداً (مثل التحكم الموجه للمجال \textenglish{Field Oriented Control - FOC})، وتكلفتها الإجمالية أعلى مقارنة بنظيرتها في التيار المستمر.
\end{itemize}

\subsection{مقارنة شاملة}
يلخص الجدول (\ref{tab:ac_dc_comparison}) أبرز الفروقات التشغيلية والهندسية بين النوعين:

\begin{table}[h!]
    \centering
    \renewcommand{\arraystretch}{1.4}
    \begin{tabular}{|p{3.5cm}|p{5cm}|p{5cm}|}
        \hline
        \textbf{وجه المقارنة} & \textbf{سيرفو التيار المستمر (\textenglish{DC})} & \textbf{سيرفو التيار المتناوب (\textenglish{AC})} \\
        \hline
        \textbf{الصيانة والعمر الافتراضي} & تتطلب صيانة (لتبديل الفرش) وعمرها الافتراضي محدود & خالية من الصيانة وعمرها الافتراضي طويل جداً \\
        \hline
        \textbf{تبديد الحرارة} & ضعيف نسبياً & ممتاز (يسهل تبريد العضو الثابت الخارجي) \\
        \hline
        \textbf{أداء السرعة والعزم} & ينخفض العزم المتاح تدريجياً مع ارتفاع السرعة & عزم ثابت ومستقر عبر نطاق واسع جداً من السرعات \\
        \hline
        \textbf{التطبيقات الشائعة} & الروبوتات المتنقلة (\textenglish{AGVs})، والمعدات الطبية، والأجهزة الصغيرة & الأذرع الروبوتية ذات الدرجات المتعددة، آلات \textenglish{CNC}، وخطوط التعبئة السريعة \\
        \hline
    \end{tabular}
    \caption[مقارنة هندسية وعملية بين محركات السيرفو للتيار المستمر والتيار المتناوب]{\centering مقارنة هندسية وعملية بين محركات السيرفو للتيار المستمر والتيار المتناوب}
    \label{tab:ac_dc_comparison}
\end{table}